\section{Estado del arte y protocolo derivado}
\label{sec:protocolo}

\subsection{Marco te\'orico y an\'alisis del estado del arte}
El protocolo de transferencia inconsciente 1-de-$N$ (\emph{1-out-of-$N$ OT})
permite que un emisor con $N$ secretos $m_1,\dots,m_N$ interact\'ue con un
receptor que posee un \'indice privado $i \in \{1,\dots,N\}$, de modo que al
finalizar la ejecuci\'on el receptor aprenda $m_i$ sin obtener informaci\'on
alguna sobre los restantes $N-1$ mensajes, mientras que el emisor no aprende
absolutamente nada sobre el valor de $i$ \cite{boneh-shoup-v06}.

En la literatura existen m\'ultiples aproximaciones para construir esquemas OT:
\begin{itemize}
\item \textbf{Construcci\'on pedag\'ogica de Boneh y Shoup (\S11.6.1)
\cite{boneh-shoup-v06}:} Propone un OT 1-de-$N$ derivado del cifrado ElGamal
en un grupo c\'iclico de orden primo $G=\langle g\rangle$. El emisor sortea
$\beta$ y publica $v = g^\beta$; el receptor camufla su \'indice calculando
$u = g^\alpha \cdot v^{-i}$ con un escalar fresco $\alpha$; el emisor deriva
$N$ claves p\'ublicas ElGamal impl\'icitas $u_j = u \cdot v^j$ y encripta cada
mensaje reutilizando $\beta$. Con las salvaguardas de sesi\'on y AEAD sugeridas
en el Remark~11.2, esta construcci\'on es limpia, rigurosa y verificable,
siendo la referencia obligatoria adoptada en este proyecto.
\item \textbf{Naor y Pinkas (2001, 2005) \cite{naor-pinkas-soda01,naor-pinkas-joc05}:}
Permiten construir un OT 1-de-$N$ a partir de $\lceil \log_2 N \rceil$
instancias paralelas de OT 1-de-2. Aunque optimizan el ancho de banda asint\'otico,
introducen m\'as rondas de negociaci\'on y mayor complejidad de estado para
el rango moderado de nuestro cat\'alogo ($N \in [500, 5000]$).
\item \textbf{Tzeng (2002) \cite{tzeng-pkc02}:} Propone un esquema directo de
una sola ronda; sin embargo, para garantizar la seguridad frente a receptores
maliciosos requiere pruebas de conocimiento cero (ZKP) sobre la construcci\'on
del mensaje del cliente, elevando la sobrecarga algor\'itmica.
\item \textbf{Chou y Orlandi (2015) \cite{chou-orlandi-simplestOT}:} Dise\~nan
un OT de alta velocidad (\emph{The Simplest OT}), pero su implementaci\'on en
curvas el\'ipticas con cofactor $h>1$ exige un manejo cuidadoso de subgrupos
para evitar fugas de bits, lo que motiva el uso de \texttt{ristretto255}.
\item \textbf{OT adaptativo y OPRF de Boneh--Shoup (\S11.6.3--\S11.6.4):}
Permiten evaluar funciones pseudoaleatorias inconscientes con costo de
comunicaci\'on $O(1)$ por consulta; no obstante, el enunciado del proyecto las
declara expresamente \emph{fuera de alcance}, exigiendo el protocolo Diffie-Hellman
lineal pedag\'ogico.
\end{itemize}

Por estas razones, se descartan las alternativas anteriores y se adopta la
construcci\'on de Boneh--Shoup \S11.6.1 con las extensiones del Remark~11.2.

\subsection{Especificaci\'on del protocolo adaptado}
Sean $G = \langle g \rangle$ un grupo de orden primo $q$ (\texttt{ristretto255}),
$\mathrm{HKDF}$ la funci\'on de derivaci\'on HKDF-SHA256 \cite{rfc5869-hkdf} y
$\mathrm{AEAD}$ el cifrado AES-256-GCM \cite{nist-gcm}. Se define la cadena de
contexto global $\mathsf{ctx}$ (ver Secci\'on~\ref{sec:formatos}).

El protocolo se ejecuta en las siguientes fases:

\begin{algorithm}[ht]
\caption{Protocolo interactivo OT 1-de-$N$ para entrega privada de $K_i$}
\label{alg:ot-protocol}
\small
\begin{algorithmic}[1]
\Require \textbf{Servidor:} claves $\{K_j\}_{j=1}^N$, contexto $\mathsf{ctx}$. \textbf{Cliente:} \'indice $i \in \{1,\dots,N\}$, $\{C_j\}_{j=1}^N$, $\mathsf{ctx}$.
\Ensure \textbf{Cliente:} registro descifrado $record_i$; el servidor no aprende $i$.
\State \textbf{Fase 0 (Offline):} $C_j \gets \mathrm{AEAD.Enc}(K_j, \mathit{nonce}_j, record_j, \mathsf{ctx}\Vert \mathrm{format\_u32}(j))$. El cliente descarga $C$ entero.
\State \textbf{Paso 1 (Inicio de sesi\'on):} Cliente invoca \code{POST /ot/session}. Servidor sortea $\beta \xleftarrow{R} \mathbb{Z}_q$, calcula $v \gets g^\beta \in G$,
\State \quad genera $sid \xleftarrow{R} \{0,1\}^{128}$, guarda $(sid, \beta)$ (TTL 60\,s) y responde $(sid, v)$.
\State \textbf{Paso 2 (Consulta OT):} Cliente sortea $\alpha \xleftarrow{R} \mathbb{Z}_q$, calcula $u \gets g^\alpha \cdot v^{-i} \in G$, env\'ia $(sid, u)$ a \code{POST /ot/query}.
\State \textbf{Paso 3 (Respuesta del emisor):} Servidor extrae y destruye $\beta$ para $sid$. Valida $u \in G$. Para $j = 1 \dots N$:
\State \quad $u_j \gets u \cdot v^j$ ($u_0 = u, u_j = u_{j-1}\cdot v$); $Z_j \gets (u_j)^\beta$;
\State \quad $k_j \gets \mathrm{HKDF}(\mathsf{ctx}, Z_j, \texttt{oblivion-ti/ot-v1}\Vert sid \Vert \mathrm{format\_u32}(j))$;
\State \quad $E_j \gets \mathrm{AEAD.Enc}(k_j, \mathit{nonce}'_j, K_j, \mathsf{ctx}\Vert sid \Vert \mathrm{format\_u32}(j))$. Responde con $\{(\mathit{nonce}'_j, E_j)\}_{j=1}^N$.
\State \textbf{Paso 4 (Descifrado local):} Cliente calcula $Z_i \gets v^\alpha = g^{\alpha\beta}$; deriva $k_i$;
\State \quad abre $K_i \gets \mathrm{AEAD.Dec}(k_i, \mathit{nonce}'_i, E_i, \mathsf{ctx}\Vert sid \Vert \mathrm{format\_u32}(i))$;
\State \quad abre $record_i \gets \mathrm{AEAD.Dec}(K_i, \mathit{nonce}_i, C_i, \mathsf{ctx}\Vert \mathrm{format\_u32}(i))$.
\end{algorithmic}
\end{algorithm}

En t\'erminos de red sobre HTTP, la interacci\'on comprende:
(1) descarga en bloque de \code{catalog.json};
(2) solicitud de sesi\'on (\code{POST /ot/session} $\to (sid, v)$); y
(3) consulta OT (\code{POST /ot/query} $\to (sid, u)$ y respuesta con $\{E_j\}$).

\begin{figure}[ht]
\centering
\small
\begin{tikzpicture}[node distance=0.8cm, auto,
  party/.style={rectangle,draw,minimum width=2.2cm,minimum height=0.7cm,align=center,font=\footnotesize},
  msg/.style={font=\scriptsize}]
\node[party] (r) {Receptor\\(Cliente $i$)};
\node[party,right=4.2cm of r] (s) {Emisor\\(Servidor)};
\draw[thick] (r.south) -- ++(0,-3.2);
\draw[thick] (s.south) -- ++(0,-3.2);

\draw[-{Stealth}] ([yshift=-0.5cm]r.south) -- node[above,msg] {\code{POST /ot/session}} ([yshift=-0.5cm]s.south);
\draw[-{Stealth}] ([yshift=-1.1cm]s.south) -- node[above,msg] {$(sid, v = g^\beta)$} ([yshift=-1.1cm]r.south);
\draw[-{Stealth}] ([yshift=-1.8cm]r.south) -- node[above,msg] {\code{POST /ot/query}: $(sid, u = g^\alpha v^{-i})$} ([yshift=-1.8cm]s.south);
\draw[-{Stealth}] ([yshift=-2.6cm]s.south) -- node[above,msg] {$\{E_j, \mathit{nonce}'_j\}_{j=1}^N$} ([yshift=-2.6cm]r.south);
\end{tikzpicture}
\caption{Diagrama de secuencia e interacci\'on de la sesi\'on OT.}
\label{fig:diagrama-secuencia-ot}
\end{figure}

\subsection{Manejo de sesiones repetidas e interrumpidas}
El dise\~no resuelve de forma integral las sesiones interrumpidas y repetidas:
\begin{itemize}
\item \textbf{Sesiones interrumpidas:} Cuando el servidor genera $(sid, \beta, v)$
en el Paso~1, almacena el par en una tabla en memoria con un tiempo l\'imite de
vida (\emph{TTL}) fijado en 60 segundos. Si el cliente interrumpe la ejecuci\'on,
se desconecta o sufre un fallo de red antes de enviar el Paso~2, la entrada
caduca y $\beta$ es purgado autom\'aticamente, impidiendo fugas de memoria o
acumulaci\'on de estados obsoletos.
\item \textbf{Sesiones de un solo uso (\emph{single-use}):} Al procesar la consulta
en el Paso~3, el servidor extrae $\beta$ y lo destruye de inmediato. Si un
adversario o cliente reintenta enviar el mismo $sid$ (\emph{replay attack}),
el servidor rechaza la petici\'on con error HTTP \code{409 Conflict} o
\code{404 Not Found}.
\end{itemize}

\subsection{Justificaci\'on rigurosa de seguridad}

\subsubsection{Privacidad del receptor (ocultamiento perfecto de $i$)}
La \'unica informaci\'on que el servidor recibe del cliente dependiente del
\'indice $i$ es el elemento de grupo $u = g^\alpha \cdot v^{-i} \in G$.
Dado que $G$ es un grupo c\'iclico de orden primo $q$ y que $\alpha$ es elegido
por el cliente de forma estrictamente uniforme en $\mathbb{Z}_q$, la aplicaci\'on
$\alpha \mapsto g^\alpha$ es un isomorfismo que distribuye $g^\alpha$ de forma
uniforme sobre $G$. Como $v^{-i}$ es un elemento fijo de $G$, la operaci\'on
$u = g^\alpha \cdot v^{-i}$ constituye una traslaci\'on biyectiva sobre $G$,
actuando como una libreta de un solo uso (\emph{one-time pad}) en el exponente.
Por lo tanto, para cualquier valor fijado de $v$ e $i$, $u$ es una variable
aleatoria uniformemente distribuida sobre $G$. La vista del servidor es
estad\'isticamente independiente de $i$, garantizando privacidad incondicional
(te\'orica de la informaci\'on).

\subsubsection{Privacidad del emisor (reducci\'on al problema CDH)}
Para demostrar que el cliente no puede aprender m\'as de una clave sim\'etrica
por ejecuci\'on, consideremos un adversario $\mathcal{A}$ capaz de derivar dos
claves distintas $K_{j_1}$ y $K_{j_2}$ con $j_1 \ne j_2$. Bajo el modelo de
or\'aculo aleatorio para $\mathrm{HKDF}$, $\mathcal{A}$ debe haber consultado
los dos secretos $Z_{j_1}$ y $Z_{j_2}$, los cuales satisfacen:
\begin{equation}
Z_{j_1} = (u \cdot v^{j_1})^\beta = u^\beta \cdot (v^\beta)^{j_1}, \qquad
Z_{j_2} = (u \cdot v^{j_2})^\beta = u^\beta \cdot (v^\beta)^{j_2}
\end{equation}
Dividiendo ambos elementos en el grupo $G$:
\begin{equation}
\frac{Z_{j_2}}{Z_{j_1}} = \frac{u^\beta \cdot (v^\beta)^{j_2}}{u^\beta \cdot (v^\beta)^{j_1}} = (v^\beta)^{j_2 - j_1}
\end{equation}
Puesto que $j_1 \ne j_2$ y $1 \le |j_2 - j_1| < N < q$, la diferencia $(j_2 - j_1)$
es invertible en $\mathbb{Z}_q$. Calculando la potencia con el inverso modular
$\Delta = (j_2 - j_1)^{-1} \pmod q$:
\begin{equation}
\left( \frac{Z_{j_2}}{Z_{j_1}} \right)^\Delta = \left( (v^\beta)^{j_2 - j_1} \right)^{(j_2 - j_1)^{-1}} = v^\beta = (g^\beta)^\beta = g^{\beta^2}
\end{equation}
Por el Teorema~10.1 y el Ejercicio~10.17 de Boneh y Shoup \cite{boneh-shoup-v06},
el problema de calcular $g^{\beta^2}$ a partir de $g^\beta$ (\emph{Squaring Diffie-Hellman})
es computacionalmente equivalente al problema Computational Diffie-Hellman (CDH).
En consecuencia, la existencia de un receptor capaz de abrir dos claves en una
misma sesi\'on permitir\'ia romper el supuesto CDH en $G$.

\subsubsection{Mitigaci\'on del ataque entre sesiones (Ejercicio 11.20 y Remark 11.2)}
Si el servidor reusara el mismo par $(\beta, v)$ en m\'ultiples consultas sin
vinculaci\'on de sesi\'on, un cliente que conoci\'o $Z_i = v^\alpha$ para la
consulta de $i$ podr\'ia enviar en una segunda consulta el valor manipulado
$\hat{u} = u \cdot v$. El servidor computar\'ia $\hat{u}_j = \hat{u}\cdot v^j = u\cdot v^{j+1} = u_{j+1}$,
lo que provocar\'ia que $\hat{Z}_{i-1} = Z_i$. El cliente podr\'ia descifrar
la encapsulaci\'on $\hat{E}_{i-1}$ utilizando el secreto que ya conoc\'ia,
aprendiendo un segundo registro sin autorizaci\'on.

Esta vulnerabilidad queda totalmente neutralizada en nuestro dise\~no:
\begin{enumerate}
\item \textbf{Escalar ef\'imero fresco:} Cada sesi\'on $sid$ genera un $\beta$
nuevo e independiente, haciendo que $v$ sea impredecible y no correlacionable.
\item \textbf{Separaci\'on de dominios en HKDF:} La derivaci\'on incluye
$sid$ e \'indice $j$ en el par\'ametro \code{info}, produciendo claves $k_j$
disjuntas aun si hubiera coincidencias algebraicas en $Z_j$.
\item \textbf{Autenticaci\'on v\'ia AEAD:} Las encapsulaciones $E_j$ incluyen
$AAD_{\mathrm{ot}, j} = \mathsf{ctx}\Vert sid\Vert \mathrm{format\_u32}(j)$.
Cualquier intento de abrir un $E_j$ de otra sesi\'on falla deterministamente
con \code{InvalidTag}.
\end{enumerate}

\subsection{Conteo exacto de operaciones y costos}
El conteo preciso de operaciones algebraicas en el grupo $G$ es:
\begin{itemize}
\item \textbf{Cliente (Receptor):} Realiza exactamente \textbf{3 exponenciaciones}
(c\'alculo de $g^\alpha$, c\'alculo de $v^{-i}$ y c\'alculo de $Z_i = v^\alpha$)
y \textbf{1 multiplicaci\'on} ($g^\alpha \cdot v^{-i}$). Su costo computacional
es $O(1)$.
\item \textbf{Servidor (Emisor):} Realiza exactamente \textbf{$N+1$ exponenciaciones}
(1 para $v = g^\beta$ y $N$ para calcular cada $Z_j = u_j^\beta$) y \textbf{$N$ multiplicaciones}
de grupo (reconstrucci\'on de $u_j = u_{j-1}\cdot v$ con $u_0 = u$). Su costo
computacional es $O(N)$.
\item \textbf{Sobrecarga de red:} El cliente transmite 48\,B en la petici\'on
($sid$ de 16\,B y $u$ de 32\,B). El servidor transmite $48 + 60N$\,bytes en la
respuesta ($sid$, $v$ y $N$ encapsulaciones de 48\,B con nonces de 12\,B).
\end{itemize}
